\dossier{151}{Muscula, dite aussi Galatea}
\trad{Ici gît Muscula, dite aussi Galatea, qui vécut deux ans, deux mois et dix-sept jours ; déposée le quinzième jour avant les calendes d’août (18 juillet), sous le consulat de Gratien Auguste pour la deuxième fois et de Probus (371), dans la paix du Christ.}
\comm{La formule \emph{in pace Chr(isti)} occupe une place inhabituelle : elle clôt le texte, après la date consulaire. Comme au \nnum{140}, les deux bases transcrivent une abréviation et l’EDB décrit un chrisme sur le support (chapitre~\ref{ch:dieu_christ}). Le double nom (\emph{quae et}) n’a pas de portée religieuse. Le texte provient, selon l’EDB, du cimetière de Commodilla ; il est conservé au Lapidario Cristiano.}

\dossier{183}{Petronia, vierge, de la neuvième région}
\trad{Déposition de Petronia le sixième jour avant les calendes de décembre (26 novembre), après le consulat de Gratien Auguste pour la troisième fois et de Flavius Equitius (375) ; elle vécut vierge environ dix-neuf ans ; elle est en paix ; (elle était) de la neuvième région, du côté du lac du \emph{Cuniculus}.}
\comm{Le mot \emph{virgo} peut désigner une jeune fille non mariée ou une vierge consacrée ; rien dans le texte ne permet de choisir (chapitre~\ref{ch:foi_eglise}). L’indication de domicile dans la neuvième région augustéenne, \emph{a lacu cuniculi}, est exceptionnelle dans le corpus ; la localisation de ce repère relève de la topographie de Rome et n’est pas proposée ici. L’EDB rattache le texte au cimetière de Zénon et décrit un chrisme, un paon et un rameau d’olivier ; la plaque est conservée au Collegio Teutonico.}

\dossier{202}{Cinnamius Opas, lecteur du titre de Fasciola}
\trad{Cinnamius Opas, lecteur du titre de Fasciola, ami des pauvres, qui vécut quarante-six ans, sept mois et neuf jours ; déposé en paix le dixième jour avant les calendes de mars (20 février), sous le consulat de Gratien pour la quatrième fois et de Merobaudes (377).}
\comm{L’un des trois lecteurs rattachés à un titre de la Ville (chapitre~\ref{ch:rangs}). Le \emph{titulus Fasciolae} est de nouveau nommé au \nnum{421} (398), avec le lecteur Pascentius. L’âge du défunt montre que la fonction de lecteur pouvait être exercée à l’âge mûr. La qualification \emph{amicus pauperum} associe le clerc à l’assistance. Le texte n’est pas apparié à l’EDB.}

\dossier{209}{Victorinus achète son emplacement au fossoyeur Martinianus}
\trad{Cet emplacement est celui de Victorinus, qui l’a acheté au fossoyeur Martinianus, sous (le consulat de) Valens pour la sixième fois et de Valentinien (pour la deuxième fois) (378).}
\comm{Le plus ancien acte d’achat daté du corpus qui nomme le vendeur, un \emph{fossor} (chapitre~\ref{ch:sepultures}). L’acte, sans commémoration, n’entre pas dans l’ensemble C. L’EDB restitue l’itération du second consul (\emph{Valentin[iano bis]}) et décrit, en tête du texte, un monogramme du Christ avec alpha et oméga : le signe chrétien marque ici un document de propriété.}

\dossier{214}{\emph{Ancilla dei} : « prie pour ce seul enfant »}
\trad{Ici repose une servante de Dieu, qui de ses biens a possédé tout entière cette demeure, que ses amies pleurent et dont elles cherchent le réconfort. Prie pour ce seul (enfant), la descendance que tu as laissée survivante. — (EDB seule :) Tu demeureras dans le repos éternel de la félicité ; le [quantième anomal] jour avant les calendes d’octobre. Cucurbitinus et Abundantius reposent ici ensemble. Sous le consulat de nos seigneurs Gratien pour la cinquième fois et Théodose, Augustes (380).}
\comm{Le seul témoignage assuré du corpus d’une demande d’intercession adressée au défunt (chapitre~\ref{ch:mort_esperance}). La défunte, \emph{ancilla dei}, est priée d’intercéder pour l’enfant qui lui survit. L’expression \emph{de sua omnia possedit domum ista} indique qu’elle a acquis la tombe sur ses propres biens. L’EDCS s’arrête à \emph{reliquisti} ; l’EDB transmet trois lignes supplémentaires, qui donnent la date et nomment deux autres défunts inhumés dans la même tombe. L’année 380 repose sur ces lignes. La plaque (69 × 128 cm) est conservée à la villa Albani ; l’EDB décrit des oiseaux tenant une couronne.}

\dossier{231}{Celerinus, enfant \emph{fidelis}}
\trad{Innocent, il s’en est allé, (enfant) des innocents, miséricordieux il fut. Qui n’a pas souffert de ton âge et versé de pieuses larmes ? En toi était attendue l’espérance de l’avenir, par toi une gloire durable, Celerinus, mon fils. Fidèle, tu reposes en paix. Il vécut un an et quatre mois ; déposé le sixième jour avant les calendes d’octobre (26 septembre), sous (le consulat d’)Antonius et de Syagrius (382).}
\comm{L’application de \emph{fidelis} à un enfant de seize mois suggère qu’il était baptisé (chapitre~\ref{ch:foi_eglise}). Le texte développe une rhétorique de la perte : l’enfant portait l’espérance et la gloire de la famille. L’EDB lit \emph{d(iebus) VI} là où l’EDCS lit \emph{d(epositus) VI} : la seconde lecture fait de \emph{VI Kal. Oct.} la date de la déposition, la première l’intègre à l’âge. La traduction suit l’EDCS.}

\dossier{232}{\emph{Intra limina sanctorum} : une sépulture recherchée}
\trad{[…]na, servante dans le Seigneur (?), […] a nourri […] les siens ; [amie des] pauvres ; [qui pour de si grands] mérites a reçu [une sépulture à l’intérieur des] seuils des saints ; […] a reçu [ce que beaucoup désir]ent et que peu obtiennent ; [sous le consulat d’Antonius et de] Syagrius (382).}
\comm{Le texte est conservé sur sa partie droite. Trois éléments intéressent l’histoire religieuse (chapitres~\ref{ch:mort_esperance} et~\ref{ch:sepultures}). (1) \emph{in dom(ino?) cultrix} : l’abréviation est développée avec un point d’interrogation par l’éditeur ; le texte n’est retenu que dans le profil large. (2) \emph{[- - - l]imina sanctorum} : la fin de l’expression est lue ; le mot qui la précédait est restitué. (3) \emph{[- - -]n(t) et rari accipiun(t)} : la seconde moitié de l’antithèse est lue. Le texte présente la sépulture près des saints comme la récompense des mérites d’une femme charitable et comme un privilège rare. Il ne dit pas quels saints, ni où ; il est d’origine incertaine.}

\dossier{244}{Proiecta : une épigramme de Damase}
\trad{(Vers) Que dire ? Ou me taire ? La douleur même m’interdit de parler. Ce tombeau retient les larmes — sache-le — des parents de Proiecta, qui avait été unie à Primus, son mari, belle de sa grâce, contente de sa seule pudeur. Hélas, trop chérie de l’amour d’une mère malheureuse ! Apprends ce qu’elle eut de plus que bien d’autres : après les premiers liens du mariage, arrachée aux yeux de Florus, son père, elle s’en est allée, désirant monter vers la lumière éthérée du ciel. Ces consolations des pleurs, Damase les offre à tous. — Elle vécut seize ans, [dix mois] et vingt-cinq jours ; déposée le troisième jour avant les calendes de janvier (30 décembre), sous le consulat de Flavius Merobaudes et de Flavius Saturninus (383).}
\comm{L’épigramme est signée par l’évêque de Rome Damase (366–384), et la formule consulaire la date de la dernière année de son épiscopat. Elle figure dans l’édition commentée de Trout (\emph{Damasus of Rome. The Epigraphic Poetry}, 2015, n\textsuperscript{o}~51 selon la concordance de l’EDCS) et dans les \emph{Carmina Latina Epigraphica} (CLE 670). Le texte ne nomme ni le Christ ni Dieu ; l’espérance de la défunte est exprimée dans le langage de l’ascension céleste. Le rapprochement avec la Proiecta nommée sur le coffret du trésor de l’Esquilin a été discuté dans la bibliographie consacrée à ce trésor ; il n’est pas repris ici. Les mots \emph{m(enses) X} sont signalés par l’EDCS comme une addition (\emph{«  »}).}

\dossier{249}{« Qui a reçu l’Esprit saint »}
\trad{[…] qui a reçu l’[Esprit] saint, [le seizième jour avant les calendes] de mai (16 avril ?), sous le consulat de Merobaudes [pour la deuxième fois] (383).}
\comm{Dossier examiné au chapitre~\ref{ch:esprit_trinite}. Sont lus \emph{-tus est qui ac-}, \emph{-tum spiritum}, \emph{Mai(as) Merobaud-}. Le verbe (\emph{ac[cepti]} dans l’EDCS, \emph{ac[cepit]} dans l’EDB), le début de l’adjectif, le quantième et l’itération du consulat sont restitués. La mention de l’Esprit saint est probable ; le rite qu’elle désigne et sa date ne peuvent être précisés. Le fragment (20 × 23,5 cm) provient, selon l’EDB, du cimetière de Nicomède, sur la via Nomentana.}

\dossier{268}{Lucifer, par son fils Benedictus}
\trad{Ce tombeau de son père, son fils Benedictus a voulu qu’il fût fait, par amour et en souvenir de son père, pour Lucifer, homme très innocent, [ami] des amis, qui, tant qu’il fut dans les choses humaines, parut très disponible envers tous et doux ; car il vécut toujours sans [faute de son] temps soixante-cinq ans, six mois et douze jours ; il mourut [en paix (?)]. — (EDB seule :) Le septième jour avant les calendes de juillet (25 juin), sous le consulat de Richomeres et de Clearchus (384).}
\comm{L’EDCS s’arrête à \emph{decess[it]} ; l’EDB transmet la dernière ligne, qui donne la date et fonde l’année 384 (chapitre~\ref{ch:fondement}). L’EDB souligne de nombreuses lettres, lues par les éditeurs anciens et aujourd’hui perdues. Le nom du défunt, \emph{Lucifer}, et celui du fils, \emph{Benedictus}, appartiennent l’un à l’onomastique traditionnelle, l’autre à l’onomastique chrétienne (chapitre~\ref{ch:noms}). La plaque (64 × 98 cm) est conservée à la Galleria Lapidaria des Musées du Vatican.}

\dossier{275}{Quirillus, Arménien de Cappadoce}
\trad{Citoyen arménien de Cappadoce, du nom de Quirillus, pieux, aimé de tous ; il vécut environ soixante-dix ans ; il nous a précédés dans la paix le sixième jour avant les ides de mars (10 mars), après le consulat de Richomeres et de Clearchus (385), un lundi ; (tombe) double avec Hyperechia, sa compagne.}
\comm{L’origine est mise en tête du texte, à l’accusatif (\emph{civem Armeniacum Cappadocem}), construction que l’EDCS signale. Le 10 mars 385 est bien un lundi (chapitre~\ref{ch:calendrier}). L’expression \emph{praecessit ad pacem} présente la mort comme un départ qui précède celui des survivants. La tombe est un \emph{bisomus} préparé pour le couple. L’EDB décrit un monogramme avec alpha et oméga ; la plaque est conservée au Lapidario Cristiano.}

\dossier{287}{Le jour, la lune et le signe}
\trad{[…] Honorius, [très noble enfant, et d’Euodius], \emph{vir clarissimus}, consuls (386), [le dixième jour avant les calendes de septembre (23 août)], un dimanche, [à la … heure de la nuit], au douzième jour de la lune, sous le signe [du Capricorne] ; […] trente jours ; […] le septième ; […] repose en paix ; […].}
\comm{Dossier examiné au chapitre~\ref{ch:calendrier}. Le quantième est entièrement restitué ; le jour de la semaine, l’âge de la lune et le signe zodiacal, en partie conservés, sont compatibles avec la restitution, mais ne la confirment pas indépendamment, puisque l’éditeur a pu la déduire de ces mêmes données. Les deux bases donnent le même texte ; l’EDB rattache la plaque au cimetière de Commodilla.}

\dossier{306}{Maximilla et Lucceia}
\trad{(Vers) Maximilla, capable [de garder l’alliance] du Christ et son propos : la mort lui a enlevé tous les biens de la vie ; mais il vit mieux, celui qui espère la vie éternelle du Christ, s’il médite en son âme la foi unique donnée à tous ; alors la mort elle-même pourra mourir, et la vie sera perpétuelle pour tout le temps ; il accordera à tous un pardon unique, mais selon les mérites. Toi, Lucceia, qui la soutiens, tu as fait (ce monument) au-dessus d’elle ; car, te souvenant d’elle quand tu acquittes tes vœux au tombeau, tu montres un esprit rempli d’une vie religieuse. — (Prose) D’une foi et d’une innocence admirables et d’un excellent propos, Maximilla, vierge, servante de Dieu, citoyenne de Pannonie, qui vécut cinquante ans, a été mise dans cette sépulture par Lucceia, \emph{clarissima femina}, fille de Viventius, \emph{vir} de mémoire clarissime, ancien préfet du prétoire et de la Ville éternelle, en raison de l’amitié qui les unissait ; elle a fini ses jours en même temps que sa mère Nunita, qui fut épouse d’un diacre et vécut soixante-dix ans ; le quatrième jour avant les nones d’octobre (4 octobre), sous le consulat des Flavii Timasius et Promotus, \emph{viri clarissimi} (389).}
\comm{Le texte le plus développé du corpus sur la vie éternelle (chapitres~\ref{ch:dieu_christ} et~\ref{ch:mort_esperance}). Il est gravé sur un sarcophage de marbre (72 × 212 cm) de l’aire à ciel ouvert de Saint-Sébastien. Le poème (lu dans les deux bases, sauf le début du premier vers) exprime une espérance construite : la vie éternelle du Christ, la mort de la mort, un pardon accordé selon les mérites. La partie en prose présente un milieu sénatorial (Lucceia, fille d’un ancien préfet) et clérical (Nunita, \emph{matrona diaconi}, selon la correction de l’EDCS). La ligne qui qualifie Maximilla de \emph{virgo ancilla dei civis Pannonia} manque dans le texte de l’EDB ; elle doit être contrôlée sur l’ICVR V 13355 avant toute exploitation. Le verbe \emph{[ve]getans} et la construction \emph{su[per h]anc} sont en partie restitués ; la traduction de la phrase qui les contient reste incertaine.}

\dossier{311}{Nisus : \emph{deposita in Chr(isto)}}
\trad{Déposé, Nisus, la veille des nones ; déposée dans le Christ le sixième jour avant les ides de février (8 février), sous le consulat de Maximus (388).}
\comm{Texte tracé dans le mortier frais, qui commémore deux défunts, dont le second n’est pas nommé dans la partie transmise. \emph{In Chr(isto)} est transcrit par les deux bases, mais l’EDB décrit un chrisme sur le support : l’abréviation développée rend probablement un monogramme (chapitre~\ref{ch:dieu_christ}). \emph{Maximo cons(ule)} désigne le consulat de Magnus Maximus, en usage à Rome au début de 388 ; l’EDB date de 389 et note que sa transcription diffère de l’ICVR. Datation contestée (chapitre~\ref{ch:fondement}).}

\dossier{321}{La colonne de Saint-Paul}
\trad{(c) Colonne de l’apôtre Paul ; jour de dédicace le [treizième] jour avant les calendes [de décembre (19 novembre)], sous le consulat de notre seigneur] Valentinien Auguste pour la quatrième fois et de Neoterius, \emph{vir clarissimus} (390), sous l’administration de Flavius Philippus, \emph{vir} [\emph{clarissimus}, …], le curateur étant Flavius Anastasius, [\emph{vir clarissimus}], tribun prétorien. — (a) Siricius, évêque, dévoué de tout son cœur. — (b) L’évêque dévoué.}
\comm{Seul texte non funéraire de l’ensemble B, inclus pour la datation consulaire et pour les noms de fonctions. L’EDB distingue trois textes gravés en trois endroits du fût : sous le chapiteau, au milieu de la colonne et la dédicace datée ; elle décrit un monogramme avec alpha et oméga après \emph{episcopus}. Le texte atteste l’engagement de l’évêque Sirice dans la construction de la nouvelle basilique de la via Ostiensis. Les noms de Flavius Philippus (fonctionnaire) et de Paul (l’apôtre) avaient été pris pour des noms consulaires par l’identification automatique (chapitre~\ref{ch:fondement}). Le quantième et une partie de la formule sont restitués.}

\dossier{324}{Hilaritas : fiançailles et mariage}
\trad{À Hilaritas, épouse bien méritante, qui vécut vingt-cinq ans ; fiancée à onze ans, elle fut fiancée sept ans ; elle vécut avec son mari sept ans et huit mois ; déposée le dix-septième jour avant les calendes de septembre (16 août), sous le consulat de Flavius Valentinien Auguste pour la quatrième fois et de Neoterius, \emph{vir clarissimus} (390).}
\comm{Le texte distingue l’âge aux fiançailles (onze ans), la durée des fiançailles (sept ans) et celle du mariage (sept ans et huit mois). Mariée vers dix-huit ans, Hilaritas est morte à vingt-cinq ans et huit mois environ ; l’âge gravé, vingt-cinq ans, est cohérent avec ce calcul (chapitre~\ref{ch:personnes}). Les deux bases concordent ; l’EDB rattache la plaque au cimetière de Commodilla.}
